% 群论（综述）
% license CCBYSA3
% type Wiki

本文根据 CC-BY-SA 协议转载翻译自维基百科\href{https://en.wikipedia.org/wiki/Group_theory}{相关文章}

在抽象代数中，群论研究的是称为“群”的代数结构。群的概念是抽象代数的核心：其他著名的代数结构，如环、域和向量空间，都可以看作是带有附加运算和公理的群。群在数学的许多领域中反复出现，而群论的方法也影响了代数的许多部分。线性代数群和李群是群论的两个分支，它们各自都有显著的发展，逐渐成为独立的研究领域。

许多物理系统，例如晶体和氢原子，以及宇宙中已知四种基本相互作用力中的三种，都可以通过对称群来建模。因此，群论及其密切相关的表示论在物理学、化学和材料科学中有着重要的应用。群论在公钥密码学中同样占据核心地位。

群论的早期历史可以追溯到19世纪。20世纪最重要的数学成就之一，是一个合作性的巨大工程，耗费了超过一万页期刊论文，主要发表于1960年至2004年，最终完成了有限单群的完整分类。
\subsection{历史}
群论有三个主要的历史来源：数论、代数方程理论和几何学。数论的脉络由莱昂哈德·欧拉开启，后由高斯在模算术以及与二次域相关的加法群和乘法群的研究中发展。关于置换群的早期成果来自拉格朗日、鲁菲尼和阿贝尔，他们在寻找高次多项式方程的一般解时做出了贡献。埃瓦里斯特·伽罗瓦创造了“群”这一术语，并建立了群论与域论之间的联系，这就是如今所称的伽罗瓦理论。在几何学中，群首先在射影几何中变得重要，后来又在非欧几何中发挥作用。费利克斯·克莱因的《埃尔兰根纲领》宣称群论是几何学的组织原则。

伽罗瓦在19世纪30年代率先使用群来判定多项式方程的可解性。阿瑟·凯莱和奥古斯丁·路易·柯西进一步推动了这一研究，创建了置换群理论。群的第二个历史来源源自几何情境。为了解释可能的几何体系（如欧几里得几何、双曲几何或射影几何），费利克斯·克莱因提出了“埃尔兰根纲领”。索福斯·李在1884年开始将群（现称李群）引入解析问题。第三，群最初是隐含使用，后来则明确应用于代数数论中。

这些早期来源的不同研究范围导致了群概念的多种表述。自约1880年起，群论开始趋于统一。从那时起，群论的影响不断扩大，推动了20世纪初抽象代数、表示论以及更多重要分支领域的诞生。有限单群的分类是20世纪中期的一项庞大工程，完成了对所有有限单群的分类。
\subsection{群的主要类别}
被研究的群类逐渐从有限置换群和一些特殊的矩阵群例子扩展到抽象群，这些抽象群可以通过生成元与关系的表示来定义。
\subsubsection{置换群}
最早得到系统研究的群类是置换群。给定一个集合$X$和一个由$X$到自身的双射（称为置换）组成的集合$G$，如果$G$在复合与取逆运算下封闭，那么$G$就是一个在$X$上作用的群。如果$X$有$n$个元素而$G$包含所有的置换，则$G$是对称$S_n$；一般而言，任意置换群 $G$都是$X$的对称群的一个子群。凯莱提出的一种早期构造展示了：任何群都可以看作一个置换群，通过左正则表示作用在其自身上（即$X = G$）。

在许多情况下，可以利用置换群在相应集合上的作用性质来研究其结构。例如，通过这种方式可以证明，当$n \geq 5$时，交错群$A_n$是单群，即它没有任何非平凡正规子群。这个事实在说明“对于次数$n \geq 5$的一般代数方程，不能用根式求解”时发挥了关键作用。
